\section{Discussion}
This study found a MALT--SS classification signal in frozen pathology foundation-model features from salivary gland slides. The retained Hybrid model combined high internal validation discrimination with external AUROC 0.8721. The more consequential finding for model development was that this signal was also accessible to substantially simpler classifiers, while controlled experiments did not establish a stable advantage for the full aggregation architecture. Together, these results support further evaluation of the disease-specific task and provide a clearer basis for choosing which parts of the pipeline require improvement.

The clinical endpoint distinguishes this work from nearby applications. Automated SS classification and focus-score assessment quantify autoimmune gland pathology, whereas orbital MALT--IgG4 classification and lymphoma subtype classification address different tissues or differential diagnoses \citep{ref10,focus2025,ref11,lymphoma2025}. Our results extend image-based investigation to the recorded MALT/SS categories in salivary gland material, with whole-slide labels rather than manually annotated training regions. They do not establish performance specifically in diagnostically equivocal biopsies. Defining that case spectrum is important: success on an archive of established diagnoses need not imply the same benefit when morphology is borderline and ancillary testing is being considered.

The baseline results place the Hybrid model in an appropriate role. Its five summaries encode complementary mathematical descriptions of the sampled features, but representational variety alone does not demonstrate added diagnostic information. Matched attention with normalization used 56.75\% fewer parameters and had an external five-seed mean AUROC close to Hybrid's. Linear mean-feature models also had competitive ranking performance. These observations are compatible with substantial task information already being present in the frozen encoder space; they do not isolate the causal contribution of the encoder, because no matched cross-encoder experiment was performed. Recent MIL work likewise emphasizes training strategy and sampling when translating foundation features into slide predictions \citep{nnmil2026}. The local feature-subspace and masking experiments here were adaptations, not full reproductions of nnMIL or ACMIL \citep{acmil2024}.

Training variation and patch-sampling variation had different practical consequences. Repeated inputs to the retained checkpoint gave a comparatively narrow external AUROC range, whereas separately trained models produced more dispersed external results despite similar internal validation scores. This pattern suggests that optimization or learned representations deserve attention alongside patch coverage. It does not identify a specific source of domain sensitivity, because initialization, sampled training inputs and optimization trajectories vary together across seeds. At the individual-slide level, input sampling still mattered: a subset of slides changed threshold classification across draws. Repeated-bag evaluation can expose this dependence, although averaging did not consistently improve the error balance.

External operating-point results also depended on the evaluated material. Removing IHC/special-stain archive tags changed sensitivity much more than the AUROC point estimate. That observation argues for reporting ranking and threshold metrics separately and for defining the intended clinical input before interpreting external performance. It cannot determine the effect of actual staining, because archive labels are not confirmed stain identities and the SS series lacks a resolved stain classification. The small external SS group further limits specificity precision. The low NPV at threshold 0.5 does not support a rule-out use, and a low Brier score in this highly imbalanced collection does not by itself establish calibration. A clinically relevant threshold would need to be selected on development data for a defined use and then evaluated on new material.

The regional outputs provide a route to pathological interrogation, rather than a demonstrated explanation of diagnosis. High- and low-margin patches are selected by relative similarity to latent prototype banks; those banks have no independent regional disease supervision. Similarly, a large gate weight is a scaling coefficient within a learned fusion network, not a causal share of the prediction. Reviewing the selected patches alongside randomly sampled regions could test whether these internal signals capture relevant infiltrate composition or glandular context. Such evidence is needed before assigning disease-specific meaning to the maps or claiming reliable lesion localization.

Several boundaries determine the interpretation of this retrospective study. Internal validation was used for checkpoint selection. The external cohort was evaluated during successive development analyses, so the later algorithm comparisons are exploratory reuse of that cohort. The current numerical analyses use slides as their unit; patient linkage, clinical reference-standard details, collection periods and confirmed material categories remain to be incorporated from the hospital records. The fixed-prediction intervals do not account for patient clustering or model-selection uncertainty. Finally, this manuscript does not report a paired reader study or an AI-assisted reading experiment. Completing the clinical metadata and evaluating a frozen pipeline in a new, representative patient cohort would address a more consequential gap than additional architecture searches on the same external slides.

In conclusion, frozen pathology features supported retrospective salivary gland MALT--SS classification, but the evidence favored careful characterization of the learned signal over a claim of aggregation superiority. Simple baselines, controlled retraining and repeated inputs made the distinction between discrimination, stability and a clinical operating point explicit. These findings provide a practical foundation for a focused clinical validation study.
